\section{Protocolo de validação do RAG}

Complementarmente ao Text-to-SQL, conduziu-se avaliação da recuperação documental sobre o Diário Oficial municipal indexado em Qdrant (9\,346 \textit{chunks}, embedding \texttt{text-embedding-3-small}).

\subsection{Conjunto de referência RAG}

O \textbf{conjunto de referência RAG ampliado} contém $N_{\text{RAG}} = 40$ perguntas estratificadas em dois estratos:

\begin{itemize}
    \item \textbf{com\_numero} ($n = 32$): decretos sorteados uniformemente do inventário do índice (seed 42), com pergunta-template contendo número e ano explícitos do ato;
    \item \textbf{tematico} ($n = 8$): perguntas curadas sobre conteúdo normativo, reaproveitadas da amostra temática inicial, sem depender exclusivamente de número na formulação.
\end{itemize}

A inclusão dos itens \textit{com\_numero} é independente do \textit{rank} de recuperação --- critério de amostra honesta. Cada item possui documento de referência validado no índice (cabeçalho do ato presente em \textit{page\_content}).

\subsection{Indicadores de recuperação}

Para cada pergunta $i$, registra-se o ranking do documento esperado entre os $k$ primeiros \textit{chunks} retornados ($k = 5$):

\begin{equation}
\text{Recall@}k = \frac{1}{N_{\text{RAG}}} \sum_{i=1}^{N_{\text{RAG}}} \mathbb{1}\!\left[\text{doc esperado} \in \text{top-}k_i\right]
\end{equation}

A \textbf{taxa de recuperação média} (MRR --- \textit{Mean Reciprocal Rank}) mede a posição média do primeiro acerto:

\begin{equation}
\text{MRR} = \frac{1}{N_{\text{RAG}}} \sum_{i=1}^{N_{\text{RAG}}} \frac{1}{\text{rank}_i}
\end{equation}

onde $\text{rank}_i$ é a posição do \textit{chunk} que contém o ato esperado, ou zero se ausente do top-$k$. Reportam-se também intervalos de confiança de Wilson (95\%) \cite{Agresti1998} para as proporções de acerto.

\subsection{Escopo e limitações do protocolo RAG}

A avaliação compara estratégias de recuperação (busca densa, deduplicação, filtro por tipo de ato e busca híbrida por número/ano) sobre o mesmo \textit{golden}. Não avalia geração da resposta final em linguagem natural nem o roteamento do orquestrador.
